\documentclass[../../main.tex]{subfiles}
\begin{document}
\section{习题}\label{section:ch19-problems}

\begin{exercise}\label{exercise:ch19-01}
证明命题\cref{proposition:ch19-微分理想判别}（一个光滑分布对合当且仅当它确定一个微分理想）。
\end{exercise}

\begin{exercise}\label{exercise:ch19-02}
设 $D$ 是光滑 $n$ 维流形 $M$ 上一个秩为 $k$ 的光滑分布，$\omega^1,\dots,\omega^{n-k}$ 是 $D$ 在开子集 $U\subseteq M$ 上的光滑局部定义形式。证明 $D$ 在 $U$ 上对合当且仅当对每个 $i=1,\dots,n-k$，下列恒等式成立：
\begin{align*}
\mathrm{d}\omega^i\wedge\omega^1\wedge\cdots\wedge\omega^{n-k}=0.
\end{align*}
（用于第 582 页。）
\end{exercise}

\begin{exercise}\label{exercise:ch19-03}
设 $\omega$ 是光滑流形 $M$ 上的一个光滑 $1$-形式。若光滑正函数 $\mu$ 在开子集 $U\subseteq M$ 上使 $\mu\omega$ 在 $U$ 上恰当，则称 $\mu$ 为 $\omega$ 在 $U$ 上的一个\emph{积分因子}。证明下列命题：
\begin{enumerate}[(a)]
\item 若 $\omega$ 处处非零，则 $\omega$ 在每一点的一个邻域内存在积分因子当且仅当 $\mathrm{d}\omega\wedge\omega\equiv 0$。
\item 若 $\dim M=2$，则每个非零光滑 $1$-形式在每一点的一个邻域内都存在积分因子。
\end{enumerate}
\end{exercise}

\begin{exercise}\label{exercise:ch19-04}
设 $U\subseteq\mathbb{R}^3$ 是三个坐标全为正的子集，$D$ 是 $U$ 上由下列向量场张成的分布：
\begin{align*}
X=y\frac{\partial}{\partial z}-z\frac{\partial}{\partial y},\qquad
Y=z\frac{\partial}{\partial x}-x\frac{\partial}{\partial z}.
\end{align*}
求 $D$ 在 $U$ 上的一个显式整体平坦坐标图。
\end{exercise}

\begin{exercise}\label{exercise:ch19-05}
设 $D$ 是 $\mathbb{R}^3$ 上由下列向量场张成的分布：
\begin{align*}
X=\frac{\partial}{\partial x}+yz\frac{\partial}{\partial z},\qquad Y=\frac{\partial}{\partial y}.
\end{align*}
\begin{enumerate}[(a)]
\item 求 $D$ 的一条过原点的积分子流形。
\item $D$ 是否对合？根据 (a) 解释你的答案。
\end{enumerate}
\end{exercise}

\begin{exercise}\label{exercise:ch19-06}
设 $D$ 是光滑流形 $M$ 上的一个对合分布，$\gamma\colon J\to M$ 是一条光滑曲线。证明下列命题：
\begin{enumerate}[(a)]
\item 若 $H$ 是 $D$ 的一个积分流形，且 $\gamma$ 的像含于 $H$，则对所有 $t\in J$ 都有 $\gamma'(t)\in T_{\gamma(t)}H\subseteq T_{\gamma(t)}M$。
\item 反之，若对所有 $t$ 都有 $\gamma'(t)\in D$，则 $\gamma$ 的像含于 $D$ 所确定的叶状结构的单个叶中。
\end{enumerate}
[注：将此与习题 5-19 的结果比较。]
\end{exercise}

\begin{exercise}\label{exercise:ch19-07}
设 $D$ 是光滑流形 $M$ 上的一个对合分布，$N$ 是 $D$ 的一个连通积分流形。证明若 $N$ 是 $M$ 的一个闭子集，则它是一个极大连通积分流形，因而是 $D$ 所确定的叶状结构的一叶。（用于第 545 页。）
\end{exercise}

\begin{exercise}\label{exercise:ch19-08}
设 $M$ 与 $N$ 是光滑流形，$F\colon M\to N$ 是光滑淹没。证明 $F$ 的非空水平集的连通分支构成 $M$ 的一个叶状结构。
\end{exercise}

\begin{exercise}\label{exercise:ch19-09}
证明命题\cref{proposition:ch19-不变分布与叶状}（不变分布与不变叶状结构）。
\end{exercise}

\begin{exercise}\label{exercise:ch19-10}
设 $M$ 与 $N$ 是光滑流形。设 $\mathscr F$ 是 $M$ 上一个余维为 $k$ 的叶状结构，$\varphi\colon N\to M$ 是光滑映射。证明：若 $\varphi$ 横截于 $\mathscr F$ 的每一叶，则在 $N$ 上存在唯一的余维-$k$ 叶状结构 $\varphi^*\mathscr F$，称为 $\mathscr F$ 的\emph{拉回}，使得 $\varphi$ 把 $\varphi^*\mathscr F$ 的每一叶映入 $\mathscr F$ 的单个叶中。
\end{exercise}

\begin{exercise}\label{exercise:ch19-11}
考虑下列关于 $u\in C^\infty(\mathbb{R}^3)$ 的偏微分方程组：
\begin{align*}
x\frac{\partial u}{\partial x}+\frac{\partial u}{\partial y}+x(y+1)\frac{\partial u}{\partial z}&=xy,\\
\frac{\partial u}{\partial x}+y\frac{\partial u}{\partial z}&=y-1.
\end{align*}
求该方程组满足 $u(0,0,z)=z$ 的解 $u$。[提示：参看例\cref{example:ch19-frobenius-example}。]
\end{exercise}

\begin{exercise}\label{exercise:ch19-12}
考虑下列关于 $u\in C^\infty(\mathbb{R}^4)$ 的偏微分方程组：
\begin{align*}
\frac{\partial u}{\partial w}+x\frac{\partial u}{\partial y}+2(w+xy)\frac{\partial u}{\partial z}&=0,\\
\frac{\partial u}{\partial x}-w\frac{\partial u}{\partial y}+2(x-wy)\frac{\partial u}{\partial z}&=0.
\end{align*}
\begin{enumerate}[(a)]
\item 证明在 $\mathbb{R}^4$ 的任意开子集上都不存在两个微分线性无关的解 $u^1,u^2$。
\item 证明在每一点的一个邻域内都存在一个微分非零的解。
\end{enumerate}
\end{exercise}

\begin{exercise}\label{exercise:ch19-13}
在下列偏微分方程组中，判断哪些在原点的一个邻域内对任意正的 $z(0,0)$（在 (c) 中即 $z(0,0)$ 与 $w(0,0)$）都有解 $z(x,y)$（在 (c) 中即 $z(x,y)$ 与 $w(x,y)$）。
\begin{enumerate}[(a)]
\item $\dfrac{\partial z}{\partial x}=z\cos y$；$\dfrac{\partial z}{\partial y}=z\log z\,\tan y$。
\item $\dfrac{\partial z}{\partial x}=e^{x-z}$；$\dfrac{\partial z}{\partial y}=x e^{y-z}$。
\item $\dfrac{\partial z}{\partial x}=z$；$\dfrac{\partial z}{\partial y}=w$；$\dfrac{\partial w}{\partial x}=w$；$\dfrac{\partial w}{\partial y}=z$。
\end{enumerate}
\end{exercise}
\end{document}
